\rok{Фебруар 2026.}{А}

% Информације о тесту 
Други практични тест одржан 10.02.2026. године. Време израде теста је 90 минута. Ово је \textit{closed book test}, није дозвољено коришћење додатних материјала. Укупан број бодова на тесту је 25. За израду се користи LINQPad шаблон који се мора преименовати у \texttt{AISP2025\_Test2\_GrupaA\_IT1g2024}.

\zadatak{1}{Конверзија листе у матрицу суседства}
Написати алгоритам који ће вршити директну конверзију листе суседства у матрицу суседства. Алгоритам као улазни параметар треба да прихвати динамички низ чворова, а треба да резултатује креираном матрицом.

\textbf{Додатне напомене за решавање задатка:}
\begin{itemize}
    \item У шаблону је закоментарисана метода \texttt{ConvertAdjacencyListToAdjacencyMatrix} која треба да садржи имплементацију конверзије листе у матрицу суседства.
    \item У \texttt{Main} методи обезбедити начин за тестирање и проверу нове функционалности.
\end{itemize}

\zadatak{2}{Проналажење првог карактера који се не понавља}
Написати алгоритам временске комплексности $O(n)$ који идентификује и исписује први карактер у прослеђеном стрингу који се у том стрингу појављује само једном.

 \textbf{Додатни захтеви за решавање задатка:}
\begin{itemize}
    \item Користити \textit{Hash} табелу из шаблона за чување броја појављивања сваког карактера.
    \item Имплементацију писати у методи: \\ \texttt{public static void firstUnique(string str)}.
    \item Игнорисати разлику између малих и великих слова.
    \item Ако такав карактер не постоји, исписати поруку да су сви карактери дупликати.
\end{itemize}

\textbf{Пример:}
\begin{itemize}
    \item \textbf{Улаз 1:} ,,Analiza''  \quad\textbf{Излаз 1:} 'n' 
    \item \textbf{Улаз 2:} ,,Baba'' \quad\textbf{Излаз 2:} ,,Svi karakteri su duplikati''
\end{itemize}

\zadatak{3}{Проналажење другог највећег елемента (BST)}
У оквиру класе \texttt{BinarySearchTree} написати методу која проналази другу по величини вредност у стаблу.

\textbf{Додатни захтеви за решавање задатка:}
\begin{itemize}
    \item Јавна метода: \texttt{public int getSecondLargest()}.
    \item Приватна метода: \texttt{private void getSecondLargest(Node current, int counter)}. (Савет: Можете користити бројач или претрагу десног подстабла).
    \item Подразумева се да стабло има барем два чвора.
    \item Алгоритам \textbf{ОБАВЕЗНО} писати у оквиру класе \texttt{BinarySearchTree}.
\end{itemize}

\noindent \textbf{Примери:}

\begin{center}
\begin{minipage}{\textwidth}
\centering
\begin{forest}
for tree={
    circle,
    draw,
    thick,
    fill=gray!20,
    minimum size=8mm,
    inner sep=1pt,
    s sep=8mm,
    l sep=12mm,
    edge={-, thick},
    font=\small
}
[8
    [4
        [2
            [, phantom]
            [3]
        ]
        [5]
    ]
    [9
        [, phantom]
        [11]
    ]
]
\end{forest}

\vspace{0.3cm} % Мали размак између стабла и текста
\textbf{Слика 1.} Пример BST-a.
\end{minipage}
\end{center}


\begin{itemize}
    \item \textbf{Улаз 1:} BST са слике 1. 
    \item \textbf{Излаз 1:} 9 
\end{itemize}




\begin{center}
\begin{minipage}{\textwidth}
\centering
\begin{forest}
for tree={
    circle,
    draw,
    thick,
    fill=gray!20,
    minimum size=8mm,
    inner sep=1pt,
    s sep=8mm,
    l sep=12mm,
    edge={-, thick},
    font=\small
}
[10
    [4
        [2
            [3]
            [, phantom]
        ]
        [6
            [5]
            [9]
        ]
    ]
    [20
        [12
            [, phantom]
            [14]
        ]
        [29
            [27]
            [, phantom]
        ]
    ]
]
\end{forest}

\vspace{0.3cm} % Мали размак између стабла и текста
\textbf{Слика 2.} Пример BST-a.
\end{minipage}
\end{center}



\begin{itemize}
    \item \textbf{Улаз 2:} BST са слике 2.
    \item \textbf{Излаз 2:} 27
\end{itemize}

